\documentclass[11pt,a4paper]{coverletter}

\senderblock
  {Radheem Bin Razi}
  {Distributed Systems \& Cloud-Native Engineer}
  {Ilmenau, Germany \textbar\ \href{mailto:sheikh.radheem@gmail.com}{sheikh.radheem@gmail.com} \textbar\ \href{https://radheem.github.io/my-cv}{portfolio}}

\begin{document}

%% ===== ENGLISH =====
\recipient{FION Energy}{Hiring Team}{}

\opening{Dear Hiring Team,}

The part of engineering I value most is the seam where software meets physical hardware, because that is where theoretical designs finally face real-world constraints. Building reliable systems that survive unpredictable networks and aging components is exactly what draws me to the IoT Platform Engineer role at FION Energy. Your mission to optimize industrial battery operations through a unified edge-to-cloud platform aligns perfectly with my background in shipping distributed systems that must operate continuously in production environments.

Over the past five years I have designed and deployed production-grade infrastructure that bridges cloud orchestration with constrained edge environments. I built a Kubernetes-based 5G testbed where I managed virtual radio interfaces, routed traffic across network namespaces, and maintained strict latency requirements for active user equipment. That experience taught me how to monitor distributed components, automate failover, and debug connectivity issues without direct physical access. I also architected event-driven microservices using Go and NATS JetStream to route data between multiple third-party vendors, which required strict schema validation and resilient error handling when upstream systems changed unexpectedly. These projects required me to treat deployment artifacts as living systems rather than static code, a mindset that translates directly to managing a growing fleet of energy storage controllers.

I am comfortable switching between high-level platform architecture and hands-on debugging, and I thrive in environments where early engineers directly shape the product roadmap. I am available to start full-time immediately, and I can relocate anywhere within Germany within two to three weeks. I am based in Germany and hold Blue Card eligibility, so your team only needs to issue the employment contract without any visa sponsorship overhead. I would welcome the chance to discuss how my experience with edge orchestration and production reliability can support FION Energy’s next phase of growth.

\closing{Sincerely,}{Radheem Bin Razi}

\clearpage
\selectlanguage{ngerman}

%% ===== DEUTSCH =====
\recipient{FION Energy}{Personalabteilung}{}

\opening{Sehr geehrtes Hiring-Team,}

Der Aspekt der Ingenieurarbeit, den ich am meisten schätze, ist die Schnittstelle, an der Software auf physische Hardware trifft, denn genau dort stoßen theoretische Entwürfe erstmals auf reale Einschränkungen. Der Aufbau zuverlässiger Systeme, die unvorhersehbare Netzwerke und alternde Komponenten überstehen, ist genau der Grund, warum mich die Rolle als IoT Platform Engineer bei FION Energy anspricht. Ihr Ziel, industrielle Batteriebetriebe durch eine einheitliche Edge-to-Cloud-Plattform zu optimieren, deckt sich perfekt mit meinem Hintergrund in der Bereitstellung verteilter Systeme, die in Produktionsumgebungen durchgehend betriebsbereit sein müssen.

In den vergangenen fünf Jahren habe ich produktionsreife Infrastruktur konzipiert und implementiert, die Cloud-Orchestrierung mit ressourcenbeschränkten Edge-Umgebungen verbindet. Ich habe ein Kubernetes-basiertes 5G-Testfeld aufgebaut, in dem ich virtuelle Funkinterfaces verwaltete, Traffic über Network Namespaces routete und strikte Latenzanforderungen für aktives User Equipment einhielt. Diese Erfahrung hat mir beigebracht, wie man verteilte Komponenten überwacht, Failover-Prozesse automatisiert und Verbindungsprobleme ohne direkten physischen Zugriff debuggt. Zudem habe ich ereignisgesteuerte Microservices mit Go und NATS JetStream architekuriert, um Daten zwischen mehreren Drittanbietern zu routen, was bei unerwarteten Änderungen in vorgelagerten Systemen eine strikte Schema-Validierung und resiliente Fehlerbehandlung erforderte. Für diese Projekte musste ich Deployment-Artefakte als lebende Systeme und nicht als statischen Code betrachten – eine Denkweise, die sich direkt auf die Verwaltung einer wachsenden Flotte von Energiespeicher-Controllern übertragen lässt.

Ich bewege mich gleichermaßen sicher zwischen hochgradiger Plattformarchitektur und praktischem Debugging und entfalte mich am besten in Umgebungen, in denen frühe Ingenieure die Produktroadmap direkt mitgestalten. Ich stehe ab sofort zur Verfügung und kann innerhalb von zwei bis drei Wochen an jeden beliebigen Ort in Deutschland umziehen. Ich bin bereits in Deutschland ansässig und erfülle die Voraussetzungen für die Blue Card, sodass Ihr Team den Arbeitsvertrag lediglich ausstellen muss, ohne zusätzlichen bürokratischen Aufwand im Rahmen eines Visa-Sponsorings. Gerne würde ich die Gelegenheit nutzen, zu besprechen, wie meine Erfahrung in Edge-Orchestrierung und Produktionszuverlässigkeit die nächste Wachstumsphase von FION Energy unterstützen kann.

\closing{Mit freundlichen Grüßen,}{Radheem Bin Razi}

\end{document}
